\section*{Discussion}

The matched experiment directly addresses whether a frozen sepsis AI system retains its behavior under consent-related input restrictions. The answer is conditional: the magnitude and operational meaning of the change depend on the model, hospital, restricted source and endpoint. All model--hospital mean AUROCs decline under the tested restrictions. Some declines are small, but a small ranking change does not establish preserved incident detection. The hospital-B GRU laboratory restriction is the clearest example: nearly unchanged AUROC coexists with a substantial increase in missed incidents at the original operating threshold.

This motivates a benchmark contribution centered on the evaluation design. Fixed patients, predictions generated by the same fitted model, explicit permission interventions and paired uncertainty allow another method to be compared on the same question. Reporting source availability, ranking, missed incidents and review burden exposes tradeoffs that an AUROC-only benchmark would conceal. The present release establishes reference baselines and a reproducible stress test; it does not establish that these baselines represent the best achievable performance under restricted data. Models trained with permission variation, source-specific architectures and validation-only adaptation are natural additional benchmark entries.

The broader research question also asks whether sensitive information is used with permission and whether consequential actions receive effective oversight. The graph representation makes specified permissions and review paths testable. The statistical construction evaluates joint risk and workload on a frozen grid. The policy evaluation produces a negative answer for the tested first-alert family: the declared incident-risk target is not achieved, despite useful discrimination and positive challenge utility. Actual clinician effectiveness and patient benefit were not measured. Risk tolerances and review budgets are technical targets, not ethical thresholds derived from patients or clinicians. Changes in simulated consent affect benchmark alert performance and modeled workload; they are not measured causal effects on mortality, treatment or wellbeing.

This is more informative than observing that a bound failed to certify a weak model. The local confidence calculations exclude every tested candidate, and the test frontiers remain far above the 5\% risk target. The calibration pools are large enough that a zero-miss candidate could pass the finite-sample correction. The principal obstacle in this experiment is the combination of score trajectories and the definition of timely first review. More calibration data alone cannot turn a high population-risk candidate into a low-risk one.

The distinction also explains why the stronger predictors do not dominate the linear model on every endpoint. AUROC rewards correct ordering of individual hourly scores. Timely first review depends on the joint trajectory: scores before the window must remain below threshold, while at least one in-window score must cross it. The joint law of the pre-window and in-window maxima is therefore central. An improvement in hourly ranking need not improve that joint event, and a model that sends an early alert can receive utility credit while failing the strict first-review criterion.

Threshold nesting is often invoked to justify risk calibration, but the decision procedure sits between the score and the loss. Alert suppression, cooldown, a one-review limit or a persistent clinical assessment can change the loss family even when the score thresholds are unchanged. The counterexample in the first-alert non-monotonicity proposition isolates this issue without relying on a particular dataset or model architecture.

A useful operational policy may allow an early assessment to remain active, request a second review after deterioration, or suppress repeated notifications for a fixed period. None of these policies is ruled out by our results. Each would need a precise definition of review cost, the duration of an assessment and the event that counts as a successful action. The finite-grid method can evaluate non-monotone policies once their patient-level losses and workloads are specified, although a workload variable above one would require a bounded-mean or count-aware interval rather than the Bernoulli interval used here.

The distinction between infeasible and undetermined should remain visible in such evaluations. A confidence method that only returns accepted candidates cannot, by itself, distinguish an inadequate predictor from an underpowered calibration set. Our lower-bound exclusion criterion supplies that additional information, while limiting the conclusion to the actual finite grid. Extending an exclusion claim to all real-valued thresholds would require additional structure or a uniform result beyond the present analysis.

The graph analysis and queue simulation answer complementary questions. The former asks whether a consequential action can be reached without traversing a review node. The latter asks whether a request receives a timely correct review under explicit assumptions. A path constraint can be satisfied even when every reviewer is overloaded or mistaken; a favorable queue result can coexist with an unmodeled bypass in a real software system.

The queue results are dominated by the upstream incident policy because a reviewer cannot repair a missing or mistimed request within the specified endpoint. This does not imply that reviewer quality is generally unimportant. It shows that improving one component cannot compensate for an event that the system never routes to that component in time. Clinical implementation evidence such as TREWS adoption \citep{trews} is needed to determine how review actually occurs, and externally validated alert-burden analyses such as \citet{wong2021} are needed to assess the workload that a hospital would experience.

Permission changes alter more than the number of visible columns. A causal summary can retain earlier measurements; observation counts can reveal the measurement process; a recurrent state can encode information from a source that is no longer permitted. Clearing these paths is part of the implemented consent intervention. The conservative recurrent reset illustrates a real design tension: strict removal of mixed historical state can also remove useful permitted context.

The population consequences depend on who opts out, which sources are restricted and when permissions change. Our independent per-group denial process is a controlled stress test, not a model of real consent behavior. It is particularly important not to read the resulting cost curves as evidence that patients should disclose a particular type of information. Dynamic consent includes communication and participant agency \citep{kaye2014}; predictive cost is only one input to that broader process.

The exchangeability statement clarifies the statistical boundary. A stable joint distribution of trajectories and permission states supports the certificate after the same frozen transformation is applied. Conditional outcome-ignorability of consent does not, by itself, make a changed deployment mechanism equivalent to the calibration mechanism. Hospital shift and consent shift can also occur together. Estimating a domain ratio from an initial feature vector cannot establish that the complete longitudinal patient, outcome and consent distribution has been correctly transported.

\section*{Limitations and next evaluation}
The warning window is deliberately strict and retrospective. An alert just before its start counts as a failure even if an earlier assessment would benefit a patient. Onset is reconstructed from the supplied label rather than independently adjudicated, and workload stops at a retrospectively defined time. The study consequently evaluates a well-defined benchmark requirement; it does not establish an optimal clinical timing rule. A prospective design should specify an observable review episode and account for the persistence of an assessment.

The model comparison is limited to three implementations and modest tuning budgets. It excludes dedicated missingness-decay models, ensembles, richer temporal features and many possible training objectives. Class weighting affects probability calibration, and no post-hoc probability calibrator is fitted. Models adapted during training to restricted inputs have not been evaluated, so the results are not a lower bound on the achievable performance of consent-constrained AI.

The sampling assumptions are also stronger than the available clinical record can verify. We treat benchmark patient files as independent units. Repeated patients across hospitals, temporal changes in care, correlated ward practices and label error could violate a deployment interpretation. Bootstrap intervals condition on fitted models and fixed splits; the three seeds reuse many patients. Group checks use a small set of recorded variables and do not establish broader fairness. The calibration-subset draws overlap, and no certified rule exists for a conditional test-risk exceedance distribution.

Consent and reviewers are simulated. There are no observed preference trajectories, clinician assessments, service-time measurements or prospective patient outcomes. Arrival and fatigue models are chosen sensitivity functions, not fitted behavioral laws. The graph is a model of information flow rather than an audit of a deployed EHR integration. Consent masks do not provide differential privacy, cybersecurity protection, deletion from previously trained model parameters or a complete account of institutional governance.

The most direct next benchmark extension is a comparison against models trained on permitted inputs, including missingness-aware temporal models and training with heterogeneous source availability. Training permission and inference permission should be varied separately. A second clinical task would test whether the observed ranking-versus-operating-point discrepancy is specific to sepsis. Real preference data would be needed to replace the independent denial mechanism with a defensible model of participant choices.

Continuing-review and cooldown policies should also be evaluated using workload units that reflect actual assessment episodes. Prospective human evaluation would then need to measure alert interpretation, correct action, response time, overrides and patient outcomes. These are required measurements, not outcomes supplied by the present simulation.

\section*{Conclusion}
Consent-related restrictions can be implemented and their predictive costs measured, but the tested sepsis systems do not preserve all aspects of performance when inputs change. The cost is source-, model- and site-dependent, and similar AUROC can coexist with substantially more missed incidents at a fixed threshold. Following hourly trajectories and clearing state at withdrawal connect the information constraint to its timing consequences. The benchmark therefore evaluates useful prediction and operating behavior together, while keeping permission enforcement, statistical guarantees and actual clinical oversight distinct. It provides evidence for comparing designs under explicit constraints; it does not decide which constraints patients should accept or how much clinical authority an AI system should receive.
